\section{Propositional logic}\label{sec:PropLogic}\index{proplogic|(}\thispagestyle{empty}

A \definend{proposition} is a statement that has a Boolean value,
either $T$ or $F$.
For instance, `$7$ is odd' and `$8^2-1=127$' are propositions,
with values $T$ and~$F$. 
In contrast, `$x$ is a perfect square' is not a proposition because for some 
$x$ it is~$T$ while for others it is not.

We can operate on propositions, including
conjoining two propositions as with `$5$ is prime and $7$ is prime',
or negating as with `it is not the case that $8$ is prime'.
The \definend{truth tables}\index{truth table} 
below define the behavior of
\definend{not} (sometimes called \definend{negation}),\index{logical operator!not}\index{negation}
\definend{and} (sometimes called \definend{conjunction}),\index{logical operator!and}\index{conjunction}
and \definend{or} (sometimes \definend{disjunction}).\index{logical operator!or}\index{disjunction}
% The first is a unary, or one input, operator while the
% other two are binary operators.
% For both tables, inputs are on the left while 
% outputs are on the right. 
\begin{center}\small % \vspace{0.5ex}
  \begin{tabular}[t]{c|c}
    \multicolumn{1}{c}{\ }    
        &\multicolumn{1}{c}{\tabulartext{not} $P$}  \\[-0.25ex]
    \multicolumn{1}{c}{$P$}      
        &\multicolumn{1}{c}{$\neg P$}  \\
    \hline
       $\false$  &$\true$  \\
       $\true$   &$\false$
  \end{tabular}
  \qquad
  \begin{tabular}[t]{cc|cc}
    \multicolumn{1}{c}{\ }    
        &\multicolumn{1}{c}{\ }    
        &\multicolumn{1}{c}{$P$ \tabulartext{and} $Q$}  
        &\multicolumn{1}{c}{$P$ \tabulartext{or} $Q$}  \\[-0.25ex]
    \multicolumn{1}{c}{$P$}      
        &\multicolumn{1}{c}{$Q$}      
        &\multicolumn{1}{c}{$P\wedge Q$}  
        &\multicolumn{1}{c}{$P\vee Q$}  \\
    \hline
       $\false$  &$\false$  &$\false$  &$\false$  \\
       $\false$  &$\true$  &$\false$  &$\true$  \\
       $\true$  &$\false$  &$\false$  &$\true$  \\
       $\true$  &$\true$  &$\true$  &$\true$  
  \end{tabular}
\end{center}
Thus, where `$7$ is odd' is $P$, and `$8$ is prime' is $Q$,
get the value of  `$7$ is odd and $8$ is prime' from the
right-hand table's third line of the third column: $\false$.
% The value of the disjunction `$7$ is odd or $8$ is prime' is $\true$.

In some fields, such as electrical engineering, the convention is to
write $0$ in place of~$F$ and $1$ in place of~$T$.
(In an electronic device these might 
stand for different voltage levels.)

The advantage of using symbols
that we can practically express more things.
For instance, if $P$ stands for `7 is odd', 
$Q$ stands for `9 is a perfect square', and 
$R$ means `11 is prime' then 
$(P\vee Q)\wedge \neg (P\vee (R\wedge Q))$
is too complex to comfortably state in everyday language.
We call that a propositional logic \definend{expression}
and denote it with a capital Roman letter, usually~$E$.

Truth tables help in working out the behavior of the complex 
statements,
by building them up from their components.
The table below shows the input/output behavior of 
$(P\vee Q)\wedge \neg (P\vee (R\wedge Q))$.
\begin{center}\small
  \begin{tabular}{ccc|ccccc}
    \multicolumn{1}{c}{$P$}
    &\multicolumn{1}{c}{$Q$}
    &\multicolumn{1}{c}{$R$}
    &\multicolumn{1}{c}{$P\vee Q$}
    &\multicolumn{1}{c}{$R\wedge Q$}
    &\multicolumn{1}{c}{$P\vee (R\wedge Q)$}
    &\multicolumn{1}{c}{$\neg(P\vee (R\wedge Q))$}
    &\multicolumn{1}{c}{\tabulartext{statement}}     \\ \hline
    $\false$  &$\false$  &$\false$  &$\false$  &$\false$  &$\false$  &$\true$  &$\false$    \\
    $\false$  &$\false$  &$\true$  &$\false$  &$\false$  &$\false$  &$\true$  &$\false$    \\
    $\false$  &$\true$  &$\false$  &$\false$  &$\false$  &$\false$  &$\true$  &$\false$    \\
    $\false$  &$\true$  &$\true$  &$\false$  &$\false$  &$\false$  &$\true$  &$\false$    \\[.5ex]
    $\true$  &$\false$  &$\false$  &$\false$  &$\false$  &$\false$  &$\true$  &$\false$    \\
    $\true$  &$\false$  &$\true$  &$\false$  &$\false$  &$\false$  &$\true$  &$\false$    \\
    $\true$  &$\true$  &$\false$  &$\false$  &$\false$  &$\false$  &$\true$  &$\false$    \\
    $\true$  &$\true$  &$\true$  &$\false$  &$\false$  &$\false$  &$\true$  &$\false$    
  \end{tabular}
\end{center}

The three `$\neg$', `$\wedge$', and `$\vee$' are 
\definend{operators}\index{operators!logicial}\index{logical operators}\index{propositional logic!operators} 
(or \definend{connectives}).\index{connectives!logicial}\index{logical connectives}
There are other operators; here are two common ones.
\begin{center}\small % \vspace{0.5ex}
  \begin{tabular}[t]{cc|cc}
    \multicolumn{1}{c}{\ }    
        &\multicolumn{1}{c}{\ }    
        &\multicolumn{1}{c}{$P$ \tabulartext{implies} $Q$}  
        &\multicolumn{1}{c}{$P$ \tabulartext{if and only if} $Q$}  \\[-0.25ex]
    \multicolumn{1}{c}{$P$}      
        &\multicolumn{1}{c}{$Q$}      
        &\multicolumn{1}{c}{$P\rightarrow Q$}  
        &\multicolumn{1}{c}{$P\leftrightarrow Q$}  \\
    \hline
       $\false$  &$\false$  &$\true$   &$\true$  \\
       $\false$  &$\true$   &$\true$   &$\false$  \\
       $\true$   &$\false$  &$\false$  &$\false$  \\
       $\true$   &$\true$   &$\true$   &$\true$  
  \end{tabular}
\end{center}

In a propositional logic expression, a single variable
such as~$P$ is an \definend{atom}.
An atom or its negation, $P$ or $\neg P$, is a \definend{literal}.
A \definend{clause} is a number of literals held together with logical 
connectives, such as
$P\vee\neg Q\vee\neg R$.

Two statements are 
\definend{equivalent}\index{equivalent propositional logic statements}
or \definend{logically equivalent} 
if they 
have the same values whenever we give the variables the same values. 
For instance, $P\rightarrow Q$ is equivalent to $\neg P\vee Q$, because
if we substitute $P=F, Q=F$ then they both give the value~$T$, 
if we substiture $P=F, Q=T$ then they also give the same value, etc.
That is, the statements are equivalent when they have the same final column
in their truth tables. 
We write $P\rightarrow Q\equiv\neg P\vee Q$

The set of formulas giving when statements are equivalent is 
\definend{Boolean algebra}.\index{Boolean algebra}\index{propositional logic!Boolean algebra}
For instance, these are the 
\definend{distributive laws}\index{distributive laws!for logic expressions}\index{propositional logic!distributive laws}
\begin{equation*}
  P\wedge(Q\vee R)\equiv(P\wedge Q)\vee (P\wedge R) 
  \qquad
  P\wedge(Q\vee R)\equiv(P\wedge Q)\vee (P\wedge R)  
\end{equation*}
and these are
\definend{DeMorgan's laws}\index{DeMorgan's laws!for logic expressions}\index{propositional logic!DeMorgan's laws}
\begin{equation*}
  \neg(P\wedge Q)\equiv \neg P\vee \neg Q 
  \qquad
  \neg(P\vee Q)\equiv \neg P\wedge \neg Q  
\end{equation*}

The three operators `$\neg$', `$\wedge$', and `$\vee$' 
form a a set that is complete in that we can reverse the activity above:~we 
can go from a table to an expression with that table. 
That is, we can produce expressions with any desired output behavior.
Here are two examples.
\begin{center} \small
  \begin{tabular}[t]{cc|c}
    \multicolumn{1}{c}{$P$}      
        &\multicolumn{1}{c}{$Q$}      
        &\multicolumn{1}{c}{$E_0$}  \\
    \hline
       $\false$  &$\false$  &$\true$  \\
       $\false$  &$\true$  &$\false$  \\
       $\true$  &$\false$  &$\false$  \\
       $\true$  &$\true$  &$\false$  
  \end{tabular}
  \hspace{3em}
  \begin{tabular}[t]{ccc|c}
    \multicolumn{1}{c}{$P$}
    &\multicolumn{1}{c}{$Q$}
    &\multicolumn{1}{c}{$R$}
    &\multicolumn{1}{c}{$E_1$}     \\ \hline
    $\false$  &$\false$  &$\false$  &$\false$    \\
    $\false$  &$\false$  &$\true$  &$\true$    \\
    $\false$  &$\true$  &$\false$  &$\true$    \\
    $\false$  &$\true$  &$\true$  &$\false$    \\[.5ex]
    $\true$  &$\false$  &$\false$  &$\true$    \\
    $\true$  &$\false$  &$\true$  &$\false$    \\
    $\true$  &$\true$  &$\false$  &$\false$    \\
    $\true$  &$\true$  &$\true$  &$\false$    
  \end{tabular}  
\end{center}
To produce $E_0$, 
focus on the row with output~$\true$. 
There, $P=\false$ and $Q=\false$ 
so the expression $E_0=\neg P\wedge \neg Q$ works. 
For $E_1$ there is more than one row with output~$\true$.
Target the second row with $\neg P\wedge \neg Q\wedge R$.
Target the third with $\neg P\wedge Q\wedge \neg R$ and
for the fifth row use~$P\wedge\neg Q\wedge\neg R$.
Finish by joining these parts with $\vee$'s.
\begin{equation*}
(\neg P\wedge \neg Q\wedge R)
  \vee(\neg P\wedge Q\wedge \neg R)
  \vee(P\wedge\neg Q\wedge\neg R)
\tag{$*$}
\end{equation*}

In a propositional logic expression, a single variable
such as~$P$ or~$Q$ is an \definend{atom}.
An atom or its negation such as $P$ or $\neg P$ is a \definend{literal}.
A \definend{clause} is a number of literals joined by connectives, 
so that $P\vee\neg Q\vee\neg R$ is a clause with three literals.
The form of ($*$) is important.
A propositional logic expression is in 
\definend{Disjunctive Normal form}\index{Disjunctive Normal form}\index{DNF}\index{Propositional logic!Disjunctive Normal form}
if is a disjunction of clauses, where each clause is a conjunction of literals,
as with ($*$) above.

There is another form that we use more often.
Statements of this form,
clauses connected by $\vee$'s, where inside each clause the statement is 
built from $\wedge$'s, 
are in \definend{Conjunctive Normal Form}.\index{Conjunctive Normal Form} 
One such expression is 
$(P_0\vee\neg P_1\vee\neg P_2)\wedge(P_1\vee P_2\vee\neg P_3)$.
Note that a clause evaluates to~\true{} if and only if at least
one of the literals evaluates to~\true.

A \definend{Boolean function}\index{Boolean function}\index{function!Boolean}\index{propositional logic!Boolean function}
has Boolean inputs, that is, they are either $T$ or~$F$, and Boolean outputs.
By the prior paragraphs, each single output Boolean function is
determined by a Boolean expression. 

% The \definend{\probname{$3$-Satisfiability} problem}\index{3-Satisfiability problem@$3$-\probname{Satisfiability}}
% is to decide the satisfiability of propositional logic expression
% where every clause has at most three literals.


\begin{exercises}
\begin{exercise}
There are sixteen binary logical operators.
Give all sixteen truth tables, and name the operator. 
\end{exercise}
\end{exercises}
\index{proplogic|)}

\endinput